\documentclass[a4paper,12pt]{article}

% Pacchetti utili
\usepackage[utf8]{inputenc}
\usepackage{amsmath, amssymb, amsthm}
\usepackage{mathtools}
\usepackage{geometry}
\geometry{margin=2.5cm}

\title{Variabili casuali continue, densità e distribuzione cumulativa}
\author{}
\date{}

\begin{document}
	
	\maketitle
	
	\section*{Distribuzioni di probabilità discrete}
	
	\subsection*{Definizioni generali}
	
	Sia $X$ una variabile casuale discreta che assume valori in un insieme numerabile $S \subset \mathbb{R}$.
	La \textit{funzione di probabilità} (o \textit{funzione di massa di probabilità}) di $X$ è indicata con
	\[
	f_X(x|\theta) = \Pr[X = x], \qquad \forall x \in S.
	\]
	
	\noindent
	In altre parole, la funzione $f_X$ assegna ad ogni possibile valore $x$ assunto dalla variabile casuale $X$ la probabilità che $X$ assuma proprio quel valore.
	
	Se un certo valore $x$ non appartiene all’insieme dei valori possibili $S$, allora
	\[
	f_X(x|\theta) = 0.
	\]
	
	In generale considereremo funzioni $f_X$ \textbf{parametriche}, cioè funzioni che dipendono da un insieme di parametri $\theta \in \Theta$. Formalmente possiamo scrivere:
	\[
	f_X : \mathbb{R} \times \Theta \longrightarrow [0,1],
	\]
	cioè una funzione che, dati un valore reale $x$ ed un valore di parametro $\theta$, restituisce un numero compreso fra $0$ e $1$ (una probabilità).
	
	\subsection*{Esempio: distribuzione binomiale}
	
	Un esempio importante di distribuzione di probabilità discreta è la \textbf{distribuzione binomiale}.  
	In questo caso la funzione di probabilità è data da:
	\[
	f_X(x|\theta) = \binom{n}{x} \, \theta^x \, (1-\theta)^{\,n-x}, 
	\qquad x = 0,1,\ldots,n,
	\]
	dove $\theta \in [0,1]$.
	
	\bigskip
	
	\noindent
	\textbf{Come leggere questa formula:}
	\begin{itemize}
		\item $\binom{n}{x}$ è il coefficiente binomiale, che indica il numero di modi in cui si possono scegliere $x$ successi su un totale di $n$ prove.
		\item $\theta^x$ è la probabilità di osservare $x$ successi (dato che $\theta$ rappresenta la probabilità di successo in una singola prova).
		\item $(1-\theta)^{n-x}$ è la probabilità di osservare i restanti $n-x$ fallimenti.
		\item Complessivamente, la formula dà la probabilità che in $n$ prove indipendenti, ciascuna con probabilità di successo $\theta$, si osservino esattamente $x$ successi.
	\end{itemize}
	
	\subsection*{Funzione di probabilità cumulata}
	
	La \textit{funzione di probabilità cumulata} di una variabile casuale discreta $X$ si indica con $F_X(x|\theta)$ ed è definita come:
	\[
	F_X(x|\theta) = \Pr[X \leq x]
	= \sum_{s \in S: \, s \leq x} f_X(s).
	\]
	
	In altre parole, $F_X(x|\theta)$ rappresenta la probabilità che $X$ assuma un valore minore o uguale a $x$.  
	Chiaramente,
	\[
	F_X : \mathbb{R} \times \Theta \longrightarrow [0,1].
	\]
	
	\subsection*{Esempio: distribuzione binomiale cumulata}
	
	Se $X \sim \text{Bin}(n,\theta)$, allora la funzione di probabilità cumulata è:
	\[
	F_X(x|\theta) = \sum_{z=0}^{x} \binom{n}{z}\, \theta^z (1-\theta)^{\,n-z}, 
	\qquad x = 0,1,\ldots,n,
	\]
	dove $\theta \in [0,1]$.
	
	\bigskip
	
	\noindent
	\textbf{Interpretazione del simbolo $z$:}
	\begin{itemize}
		\item La lettera $z$ è l'\textit{indice di somma}.  
		\item La funzione cumulata richiede la probabilità che $X \leq x$, quindi bisogna sommare tutte le probabilità degli eventi $\{X=0\}, \{X=1\}, \ldots, \{X=x\}$.
		\item Per questo compare la sommatoria da $z=0$ fino a $z=x$: ad ogni valore di $z$ corrisponde la probabilità di osservare esattamente $z$ successi.
		\item In pratica, $z$ è una variabile “fittizia” che serve a scorrere tutti i possibili valori fino a $x$. Non è un parametro né una variabile casuale “nuova”, ma solo l’indice che ci permette di sommare le probabilità.
	\end{itemize} 
	
	\section*{Variabile casuale continua e funzione di densità}
	
	Sia $X$ una variabile casuale continua distribuita secondo la 
	\textit{funzione di densità} $f_X(x|\theta)$:
	\[
	f_X : \mathbb{R} \times \Theta \to \mathbb{R}^+
	\]
	
	$f_X$ è una funzione di densità parametrica con parametri $\theta \in \Theta$.
	
$f_X(x|\theta)$: densità di probabilità della variabile $X$ dato un insieme di parametri $\theta$. 

$\Theta$ rappresenta lo spazio dei parametri 


$f_X$ fornisce la densità di X dati i parametri in $\theta$, non la probabilità. per ottenerla va integrata l'area sotto la curva



La densità $f_X$ descrive come la probabilità cambia lungo i valori reali. 

La funzione di ripartizione (la probabilità comulata) accumula quella probabilità fino a un certo valore $x$ 


\section*{Cos'è la densità di probabilità?}
	
	Per una variabile discreta (ad esempio il lancio di un dado) si assegna
	direttamente la probabilità a ciascun valore, ad esempio
	\[
	P(X=3) = \tfrac{1}{6}.
	\]
	
	Per una variabile continua (ad esempio l'altezza di una persona), la probabilità
	che $X$ assuma esattamente un valore $x$ è nulla. 
	Perciò non ha senso dire $P(X=1.75)$.
	
	Al posto di una probabilità puntuale, si usa la \textbf{densità di probabilità}.
	
	\subsection*{Definizione}
	La \textbf{densità di probabilità} $f_X(x)$ è una funzione tale che
	\[
	P(a \leq X \leq b) = \int_a^b f_X(x)\,dx.
	\]
	
	\subsection*{Proprietà}
	\begin{itemize}
		\item $f_X(x) \geq 0 \quad \forall x$ (non negatività);
		\item $\int_{-\infty}^{+\infty} f_X(x)\,dx = 1$ (normalizzazione);
		\item $f_X(x)$ non è una probabilità diretta: può anche assumere valori $>1$.
	\end{itemize}
	
	\subsection*{Esempio: distribuzione normale}
	La densità della variabile casuale normale standard è
	\[
	f_X(x) = \frac{1}{\sqrt{2\pi}} e^{-x^2/2}.
	\]
	
	\begin{itemize}
		\item Ha la tipica forma a campana.
		\item L'area sotto tutta la curva è pari a 1.
		\item La probabilità che $X$ cada tra $-1$ e $1$ è
		\[
		P(-1 \leq X \leq 1) = \int_{-1}^1 f_X(x)\,dx \approx 0.68.
		\]
	\end{itemize}
	
	\section*{Funzione di distribuzione cumulativa (CDF)}
	
	La \textbf{funzione di distribuzione cumulativa} (o \textit{CDF}, Cumulative Distribution Function) è definita come
	\[
	F_X(x|\theta) = \Pr[X \leq x] = \int_{-\infty}^x f_X(z|\theta)\,dz.
	\]
	
	\subsection*{Interpretazione intuitiva}
	La CDF restituisce la probabilità che la variabile casuale $X$ assuma un valore 
	minore o uguale a $x$. In pratica, rappresenta l'``area sotto la curva'' della densità $f_X$ fino al punto $x$.
	
	\subsection*{Proprietà della CDF}
	\begin{itemize}
		\item $F_X(x)$ è una funzione \textbf{crescente} (non decrescente);
		\item $0 \leq F_X(x) \leq 1$ per ogni $x$;
		\item $\lim_{x \to -\infty} F_X(x) = 0$;
		\item $\lim_{x \to +\infty} F_X(x) = 1$;
		\item $F_X$ è continua se $X$ è continua.
	\end{itemize}
	
	\subsection*{Esempio: distribuzione normale}
	Per una variabile normale standard $X \sim \mathcal{N}(0,1)$, la CDF è
	\[
	F_X(x) = \int_{-\infty}^x \frac{1}{\sqrt{2\pi}} e^{-z^2/2}\,dz.
	\]
	
	Questa funzione non ha una formula chiusa semplice, ma è tabulata e disponibile in tutti i software statistici.
	
\section*{Funzione quantile}

La \textbf{funzione quantile} è l'inversa della funzione di distribuzione cumulativa $F_X$.  

\begin{itemize}
	\item La CDF $F_X(x|\theta)$ prende un valore $x$ e restituisce una probabilità $p \in [0,1]$.  
	\item La funzione quantile $q_X(p|\theta)$ fa il contrario: prende una probabilità $p$ e restituisce il valore $x$ corrispondente.  
\end{itemize}

Formalmente:
\[
q_X(p|\theta) = \inf \{ x \in \mathbb{R} : F_X(x|\theta) \geq p \}.
\]

\subsection*{Interpretazione}
\begin{itemize}
	\item Se $p=0.5$, allora $q_X(0.5|\theta)$ è la \textbf{mediana} della distribuzione.  
	\item Se $p=0.25$, otteniamo il \textbf{primo quartile}.  
	\item Se $p=0.975$, otteniamo il \textbf{quantile al 97.5\%}, che si usa ad esempio negli intervalli di confidenza.  
\end{itemize}

\subsection*{Proprietà}
\begin{itemize}
	\item La funzione quantile mappa l'intervallo delle probabilità $[0,1]$ nei possibili valori reali della variabile casuale:
	\[
	q_X : [0,1] \times \Theta \to \mathbb{R}.
	\]
	\item È l'inversa della CDF, quindi
	\[
	q_X(p|\theta) = F_X^{-1}(p|\theta)
	\]
	se e solo se $F_X$ è strettamente crescente e invertibile.
\end{itemize}

\subsection*{Esempio pratico}
Per una variabile normale standard $X \sim \mathcal{N}(0,1)$:
\begin{itemize}
	\item $q_X(0.5) = 0$ (mediana).  
	\item $q_X(0.975) \approx 1.96$, valore che si usa nella statistica inferenziale per gli intervalli di confidenza al 95\%.  
\end{itemize}

	
\end{document}
